
% Learning objectives
\begin{frame}{Learning Objectives}
By the end of this chapter, you will understand:
\begin{itemize}
\item Fifth-Order Systems (5-OS) and the Harvard abstract model
\item Historical development from Turing machines to Harvard Mark I
\item Differences between Harvard and Von Neumann architectures
\item RISC principles and their advantages
\item Cache memory hierarchy and performance implications
\item Interpreting vs. executing instruction paradigms
\item Modern implementations and design considerations
\end{itemize}
\end{frame}

% subsection 1: Introduction to Fifth-Order Systems
\subsection{Introduction to Fifth-Order Systems}

\begin{frame}{From Fourth to Fifth-Order Systems}
\begin{columns}
\begin{column}{0.6\textwidth}
\concept{System Evolution Recap}:
\begin{itemize}
\item \textbf{0-OS}: Combinational circuits (no feedback)
\item \textbf{1-OS}: Memory circuits (first feedback loop)
\item \textbf{2-OS}: Automata (second feedback loop)
\item \textbf{3-OS}: Elementary processors (third feedback loop)
\item \textbf{4-OS}: Von Neumann computers (fourth feedback loop)
\item \textbf{5-OS}: Harvard computers (fifth feedback loop)
\end{itemize}

\vspace{0.3cm}
\concept{Fifth Loop Characteristics}:
\begin{itemize}
\item Processor + Program Memory = 4-OS
\item Adding Data Memory creates 5-OS
\item Separate memory spaces for instructions and data
\item Independent access paths enable parallel operations
\end{itemize}

\vspace{0.3cm}
\concept{Key Innovation}:
\begin{itemize}
\item Dual memory architecture
\item Simultaneous instruction and data access
\item Enhanced performance through parallelism
\end{itemize}
\end{column}
\begin{column}{0.4\textwidth}
\begin{center}
\textbf{System Order Progression}
\vspace{0.3cm}

\tikz{
% System evolution diagram
\foreach \y/\label/\desc in {
  4/5-OS/Harvard,
  3.4/4-OS/Von Neumann,
  2.8/3-OS/Processors,
  2.2/2-OS/Automata,
  1.6/1-OS/Memory,
  1/0-OS/Combinational
} {
  \draw[thick] (0,\y) rectangle (2.5,\y+0.4);
  \node at (0.7,\y+0.2) {\label};
  \node at (1.8,\y+0.2) {\desc};
}

% Highlight 5-OS
\draw[thick, red] (0,4) rectangle (2.5,4.4);

% Arrows showing progression
\foreach \y in {1.4, 2, 2.6, 3.2, 3.8} {
  \draw[->, thick] (1.25,\y) -- (1.25,\y+0.2);
}

\node at (1.25,0.5) {\textbf{Feedback Loops}};
}
\end{center}
\end{column}
\end{columns}
\end{frame}

\begin{frame}{Theoretical Foundation: Universal Turing Machine}
\begin{columns}
\begin{column}{0.6\textwidth}
\concept{Turing Machine Inspiration}:
\begin{itemize}
\item Alan Turing divided UTM tape into two parts:
\begin{itemize}
\item Machine description $M(T)$ (program)
\item Associated tape $T$ (data)
\end{itemize}
\item Two-memory concept: programs and data
\item Multi-tape TM: multiple memories, multiple heads
\end{itemize}

\vspace{0.3cm}
\concept{UTM Processing Model}:
\begin{itemize}
\item Collaboration between program and finite automaton
\item Data changes through:
\begin{enumerate}
\item Readings from program part
\item State transitions of finite automaton
\end{enumerate}
\item Complexity factors:
\begin{itemize}
\item Program size (flexible)
\item Automaton size (fixed in hardware)
\end{itemize}
\end{itemize}

\vspace{0.3cm}
\concept{Design Implications}:
\begin{itemize}
\item Minimize hardware complexity
\item Maximize program flexibility
\item Separate concerns: control vs. data
\end{itemize}
\end{column}
\begin{column}{0.4\textwidth}
\begin{center}
\textbf{UTM Tape Organization}
\vspace{0.3cm}

\tikz{
% UTM tape representation
\draw[thick] (0,3) rectangle (3,3.5);
\node at (1.5,3.25) {Program $M(T)$};

\draw[thick] (0,2.3) rectangle (3,2.8);
\node at (1.5,2.55) {Data $T$};

% Read/write head
\draw[thick] (1.4,2.8) -- (1.4,3);
\draw[thick] (1.6,2.8) -- (1.6,3);
\node at (1.5,3.7) {R/W Head};

% Finite automaton
\draw[thick, circle] (1.5,1.5) circle (0.4);
\node at (1.5,1.5) {FA};

% Arrows
\draw[<->] (1.5,1.9) -- (1.5,2.3);
\node at (2.2,2.1) {Control};

% Labels
\node at (0.2,3.25) {Instructions};
\node at (0.2,2.55) {Variables};
}
\end{center}
\end{column}
\end{columns}
\end{frame}

% subsection 2: Historical Context and Development
\subsection{Historical Context and Development}

\begin{frame}{Historical Timeline: Harvard Mark I (1944)}
\begin{columns}
\begin{column}{0.6\textwidth}
\concept{Key Historical Facts}:
\begin{itemize}
\item \textbf{March 29, 1944}: John von Neumann runs first program on Harvard Mark I
\item \textbf{Architect}: Howard Aiken (Harvard University)
\item \textbf{Manufacturer}: IBM (Endicott plant, New York)
\item \textbf{Concept Origin}: November 1937 (Aiken's proposal to IBM)
\item \textbf{Funding Approval}: February 1939 (Thomas Watson Sr.)
\end{itemize}

\vspace{0.3cm}
\concept{Design Innovation}:
\begin{itemize}
\item Two separate memories: programs and data
\item Relay-based electromechanical computer
\item General-purpose computational capability
\item Inspired by need to solve differential equations
\end{itemize}

\vspace{0.3cm}
\concept{Historical Significance}:
\begin{itemize}
\item First implementation of dual-memory concept
\item Foundation for Harvard abstract model
\item Influenced modern computer architecture
\item Demonstrated separation of concerns principle
\end{itemize}
\end{column}
\begin{column}{0.4\textwidth}
\begin{center}
\textbf{Harvard Mark I Timeline}
\vspace{0.3cm}

\tikz{
% Timeline
\draw[thick] (0,0) -- (0,4);

% Events
\foreach \y/\year/\event in {
  3.5/1937/Aiken Proposal,
  3/1939/IBM Approval,
  2.5/1944/Delivery,
  2/1944/First Program,
  1.5/1970s/Formal Recognition
} {
  \draw[thick] (-0.1,\y) -- (0.1,\y);
  \node[right] at (0.2,\y) {\textbf{\year}};
  \node[right] at (0.2,\y-0.2) {\event};
}

% Key figures
\node at (1.5,0.8) {\textbf{Key People}};
\node at (1.5,0.5) {Howard Aiken};
\node at (1.5,0.2) {Thomas Watson Sr.};
\node at (1.5,-0.1) {John von Neumann};
}
\end{center}
\end{column}
\end{columns}
\end{frame}

\begin{frame}{Von Neumann's "First Draft" (1945)}
\begin{columns}
\begin{column}{0.6\textwidth}
\concept{Document Background}:
\begin{itemize}
\item \textbf{Title}: "First Draft of a Report on the EDVAC"
\item \textbf{Author}: John von Neumann (sole credited author)
\item \textbf{Context}: EDVAC development at Moore School
\item \textbf{Distribution}: 24 typed copies by Herman Goldstine
\end{itemize}

\vspace{0.3cm}
\concept{Key Contributors}:
\begin{itemize}
\item \textbf{Herman Goldstine}: Typed and distributed document
\item \textbf{John Mauchly}: ENIAC/EDVAC development
\item \textbf{J. Presper Eckert}: Hardware design expertise
\item \textbf{Arthur Burks}: Theoretical contributions
\item \textbf{Alan Turing}: Theoretical foundation (UTM)
\end{itemize}

\vspace{0.3cm}
\concept{Document Characteristics}:
\begin{itemize}
\item Handwritten by von Neumann
\item No bibliographical references (unfinished)
\item Outlined stored-program computer principles
\item Shared memory for instructions and data
\end{itemize}
\end{column}
\begin{column}{0.4\textwidth}
\begin{center}
\textbf{EDVAC Development Team}
\vspace{0.3cm}

\tikz{
% Central node
\node[draw, circle, minimum width=1.2cm] (edvac) at (0,0) {EDVAC};

% Contributors
\node[draw, rectangle] (vn) at (0,1.8) {von Neumann};
\node[draw, rectangle] (gold) at (-1.5,1) {Goldstine};
\node[draw, rectangle] (mauch) at (1.5,1) {Mauchly};
\node[draw, rectangle] (eck) at (-1.5,-1) {Eckert};
\node[draw, rectangle] (burks) at (1.5,-1) {Burks};
\node[draw, rectangle] (tur) at (0,-1.8) {Turing};

% Connections
\draw[->] (vn) -- (edvac);
\draw[->] (gold) -- (edvac);
\draw[->] (mauch) -- (edvac);
\draw[->] (eck) -- (edvac);
\draw[->] (burks) -- (edvac);
\draw[->] (tur) -- (edvac);

% Labels
\node at (0,2.5) {\textbf{Collaborative Effort}};
\node at (0,-2.5) {\textbf{Moore School, 1945}};
}
\end{center}
\end{column}
\end{columns}
\end{frame}

% subsection 3: Harvard vs Von Neumann Models
\subsection{Harvard vs Von Neumann Models}

\begin{frame}{Architectural Comparison}
\begin{columns}
\begin{column}{0.6\textwidth}
\concept{Harvard Abstract Model}:
\begin{itemize}
\item \textbf{Memory Organization}: Separate program and data memories
\item \textbf{Access Pattern}: Simultaneous instruction and data access
\item \textbf{Bus Structure}: Independent instruction and data buses
\item \textbf{Performance}: Higher throughput, no memory conflicts
\item \textbf{Complexity}: More complex memory system
\end{itemize}

\vspace{0.3cm}
\concept{Von Neumann Abstract Model}:
\begin{itemize}
\item \textbf{Memory Organization}: Unified memory for programs and data
\item \textbf{Access Pattern}: Sequential access to instructions and data
\item \textbf{Bus Structure}: Shared bus for all memory access
\item \textbf{Performance}: Memory bottleneck (von Neumann bottleneck)
\item \textbf{Complexity}: Simpler memory system design
\end{itemize}

\vspace{0.3cm}
\concept{Terminology Clarification}:
\begin{itemize}
\item More accurate: "Abstract Models" not "Architectures"
\item Architecture has specific meaning (Amdahl, 1964)
\item Abstract models bridge theory and implementation
\end{itemize}
\end{column}
\begin{column}{0.4\textwidth}
\begin{center}
\textbf{Model Comparison}
\vspace{0.3cm}

% Harvard Model
\tikz{
\node at (1.5,3.5) {\textbf{Harvard Model}};

% Processor
\draw[thick] (1,2.8) rectangle (2,3.2);
\node at (1.5,3) {CPU};

% Program memory
\draw[thick] (0,2.2) rectangle (1,2.6);
\node at (0.5,2.4) {Prog};

% Data memory
\draw[thick] (2,2.2) rectangle (3,2.6);
\node at (2.5,2.4) {Data};

% Buses
\draw[thick] (1,2.9) -- (0.5,2.6);
\draw[thick] (2,2.9) -- (2.5,2.6);

\node at (0.3,2.8) {I-Bus};
\node at (2.7,2.8) {D-Bus};
}

\vspace{0.5cm}

% Von Neumann Model
\tikz{
\node at (1.5,2) {\textbf{Von Neumann Model}};

% Processor
\draw[thick] (1,1.3) rectangle (2,1.7);
\node at (1.5,1.5) {CPU};

% Unified memory
\draw[thick] (1,0.7) rectangle (2,1.1);
\node at (1.5,0.9) {Memory};

% Shared bus
\draw[thick] (1.5,1.3) -- (1.5,1.1);
\node at (1.8,1.2) {Bus};
}
\end{center}
\end{column}
\end{columns}
\end{frame}

\begin{frame}{Interpreting vs. Executing Instructions}
\begin{columns}
\begin{column}{0.6\textwidth}
\concept{Two Paradigms for Instruction Processing}:

\vspace{0.3cm}
\concept{Interpreting Instructions}:
\begin{itemize}
\item Instructions processed by control automaton
\item Sequence of micro-instructions generated
\item Complex control logic with many states
\item Flexible but slower execution
\item Higher hardware complexity for control
\end{itemize}

\vspace{0.3cm}
\concept{Executing Instructions}:
\begin{itemize}
\item Instructions executed in single clock cycle
\item Instruction fields directly control components
\item Zero-state automaton (no control states)
\item Faster execution, simpler control
\item UTM with zero internal states concept
\end{itemize}

\vspace{0.3cm}
\concept{Design Philosophy}:
\begin{itemize}
\item Minimize control complexity
\item Maximize execution efficiency
\item Separate instruction and data paths
\item Enable parallel processing
\end{itemize}
\end{column}
\begin{column}{0.4\textwidth}
\begin{center}
\textbf{Execution Models}
\vspace{0.3cm}

% Interpreting model
\tikz{
\node at (1.5,3.5) {\textbf{Interpreting}};

\draw[thick] (0.5,3) rectangle (2.5,3.2);
\node at (1.5,3.1) {Instruction};

\draw[->] (1.5,3) -- (1.5,2.7);

\draw[thick] (0.5,2.4) rectangle (2.5,2.7);
\node at (1.5,2.55) {Control Automaton};

\draw[->] (1.5,2.4) -- (1.5,2.1);

\draw[thick] (0.5,1.8) rectangle (2.5,2.1);
\node at (1.5,1.95) {Micro-instructions};

\node at (1.5,1.5) {Multiple cycles};
}

\vspace{0.5cm}

% Executing model
\tikz{
\node at (1.5,2.2) {\textbf{Executing}};

\draw[thick] (0.5,1.7) rectangle (2.5,1.9);
\node at (1.5,1.8) {Instruction};

\draw[->] (1.5,1.7) -- (1.5,1.4);

\draw[thick] (0.5,1.1) rectangle (2.5,1.4);
\node at (1.5,1.25) {Direct Control};

\node at (1.5,0.8) {Single cycle};
}
\end{center}
\end{column}
\end{columns}
\end{frame}

% subsection 4: RISC Architecture Principles
\subsection{RISC Architecture Principles}

\begin{frame}{Reduced Instruction Set Computer (RISC)}
\begin{columns}
\begin{column}{0.6\textwidth}
\concept{RISC Philosophy}:
\begin{itemize}
\item Keep only simplest and most frequently used instructions
\item Reject complex instructions that can be implemented as sequences
\item Simple hardware design outperforms complex designs
\item Machine code only marginally larger than CISC
\end{itemize}

\vspace{0.3cm}
\concept{Performance Benefits}:
\begin{itemize}
\item Faster execution due to simpler operations
\item Better pipeline efficiency
\item Reduced hardware complexity
\item Easier verification and testing
\end{itemize}

\vspace{0.3cm}
\concept{Design Principles}:
\begin{itemize}
\item \textbf{Simplicity}: Simple operations, simple addressing
\item \textbf{Regularity}: Consistent instruction format
\item \textbf{Efficiency}: Optimize common case
\item \textbf{Parallelism}: Enable instruction-level parallelism
\end{itemize}

\vspace{0.3cm}
\concept{Harvard + RISC Synergy}:
\begin{itemize}
\item Separate memories enable simultaneous access
\item Simple instructions reduce control complexity
\item Perfect match for high-performance systems
\end{itemize}
\end{column}
\begin{column}{0.4\textwidth}
\begin{center}
\textbf{RISC vs CISC}
\vspace{0.3cm}

\tikz{
% RISC characteristics
\node at (1.5,3.5) {\textbf{RISC}};
\draw[thick, fill=green!20] (0,3) rectangle (3,3.3);
\node at (1.5,3.15) {Simple Instructions};
\draw[thick, fill=green!20] (0,2.7) rectangle (3,3);
\node at (1.5,2.85) {Fast Execution};
\draw[thick, fill=green!20] (0,2.4) rectangle (3,2.7);
\node at (1.5,2.55) {Simple Hardware};

% CISC characteristics
\node at (1.5,2) {\textbf{CISC}};
\draw[thick, fill=red!20] (0,1.5) rectangle (3,1.8);
\node at (1.5,1.65) {Complex Instructions};
\draw[thick, fill=red!20] (0,1.2) rectangle (3,1.5);
\node at (1.5,1.35) {Variable Timing};
\draw[thick, fill=red!20] (0,0.9) rectangle (3,1.2);
\node at (1.5,1.05) {Complex Hardware};

% Performance arrow
\draw[->, thick, green] (3.2,2.85) -- (3.2,1.35);
\node at (3.7,2.1) {Performance};
\node at (3.7,1.8) {Advantage};
}
\end{center}
\end{column}
\end{columns}
\end{frame}

\begin{frame}{Hardware-Software Segregation}
\begin{columns}
\begin{column}{0.6\textwidth}
\concept{Clear Segregation Principle}:
\begin{itemize}
\item \textbf{Simple Hardware}: Built from simple circuits
\item \textbf{Complex Software}: Handles algorithmic complexity
\item \textbf{Separation of Concerns}: Hardware does what it does best
\item \textbf{Optimization}: Each layer optimized independently
\end{itemize}

\vspace{0.3cm}
\concept{Hardware Simplicity}:
\begin{itemize}
\item Most components are simple circuits
\item Only decoder subsection has moderate complexity
\item Size comparable to description size
\item Easier specification, design, verification, testing
\end{itemize}

\vspace{0.3cm}
\concept{Software Complexity}:
\begin{itemize}
\item Algorithms implemented in software
\item Compiler optimizations handle complexity
\item Flexible and updatable
\item No hardware redesign needed for new algorithms
\end{itemize}

\vspace{0.3cm}
\concept{Benefits}:
\begin{itemize}
\item Reduced development time and cost
\item Better reliability and testability
\item Easier maintenance and updates
\item Higher performance through specialization
\end{itemize}
\end{column}
\begin{column}{0.4\textwidth}
\begin{center}
\textbf{Segregation Model}
\vspace{0.3cm}

\tikz{
% Software layer
\draw[thick, fill=blue!20] (0,3) rectangle (3,3.5);
\node at (1.5,3.25) {\textbf{Complex Software}};
\node at (1.5,2.8) {Algorithms, Optimizations};

% Interface
\draw[thick, dashed] (0,2.5) -- (3,2.5);
\node at (1.5,2.6) {ISA Interface};

% Hardware layer
\draw[thick, fill=green!20] (0,1.5) rectangle (3,2.4);
\node at (1.5,1.95) {\textbf{Simple Hardware}};
\node at (1.5,1.7) {Basic Operations};

% Complexity indicators
\node at (3.5,3.25) {High};
\node at (3.5,1.95) {Low};
\draw[<->, thick] (3.3,1.7) -- (3.3,3.3);
\node at (3.8,2.5) {Complexity};

% Benefits
\node at (1.5,1.2) {\textbf{Benefits}};
\node at (1.5,0.9) {• Fast hardware};
\node at (1.5,0.6) {• Flexible software};
\node at (1.5,0.3) {• Easy verification};
}
\end{center}
\end{column}
\end{columns}
\end{frame}

% subsection 5: Cache Memory Hierarchy
\subsection{Cache Memory Hierarchy}

\begin{frame}{Cache Memory Levels}
\begin{columns}
\begin{column}{0.6\textwidth}
\concept{Three-Level Cache Hierarchy}:
\begin{itemize}
\item \textbf{L1 Cache}: Smallest, fastest (8KB-128KB)
\item \textbf{L2 Cache}: Medium size, moderate speed (1MB-8MB)
\item \textbf{L3 Cache}: Largest, slower but faster than main memory (8MB-64MB+)
\end{itemize}

\vspace{0.3cm}
\concept{L1 Cache Characteristics}:
\begin{itemize}
\item On-chip, same die as processor
\item Operates at processor speed
\item Harvard split: L1i (instructions) + L1d (data)
\item Provides fastest access to critical data
\end{itemize}

\vspace{0.3cm}
\concept{Cache Hierarchy Principles}:
\begin{itemize}
\item \textbf{Locality Principle}: Temporal and spatial locality
\item \textbf{Memory Hierarchy}: Speed vs. size trade-offs
\item \textbf{Performance}: Bridge processor-memory speed gap
\item \textbf{Cost-Effectiveness}: Balance performance and cost
\end{itemize}

\vspace{0.3cm}
\concept{Harvard Advantage}:
\begin{itemize}
\item Separate L1i and L1d eliminate conflicts
\item Simultaneous instruction and data access
\item Better cache utilization
\end{itemize}
\end{column}
\begin{column}{0.4\textwidth}
\begin{center}
\textbf{Memory Hierarchy}
\vspace{0.3cm}

\tikz{
% Processor
\draw[thick] (1,3.5) rectangle (2,3.8);
\node at (1.5,3.65) {CPU};

% L1 Cache (split)
\draw[thick, fill=green!30] (0.5,3) rectangle (1.2,3.3);
\node at (0.85,3.15) {L1i};
\draw[thick, fill=green!30] (1.8,3) rectangle (2.5,3.3);
\node at (2.15,3.15) {L1d};

% L2 Cache
\draw[thick, fill=blue!30] (0.8,2.5) rectangle (2.2,2.8);
\node at (1.5,2.65) {L2 Cache};

% L3 Cache
\draw[thick, fill=yellow!30] (0.6,2) rectangle (2.4,2.3);
\node at (1.5,2.15) {L3 Cache};

% Main Memory
\draw[thick, fill=red!30] (0.4,1.5) rectangle (2.6,1.8);
\node at (1.5,1.65) {Main Memory};

% Connections
\draw[<->] (1.2,3.5) -- (0.85,3.3);
\draw[<->] (1.8,3.5) -- (2.15,3.3);
\draw[<->] (1.5,3) -- (1.5,2.8);
\draw[<->] (1.5,2.5) -- (1.5,2.3);
\draw[<->] (1.5,2) -- (1.5,1.8);

% Speed/Size indicators
\node at (3,3.15) {Fast/Small};
\node at (3,1.65) {Slow/Large};
\draw[<->, thick] (2.8,1.8) -- (2.8,3.3);
}
\end{center}
\end{column}
\end{columns}
\end{frame}

\begin{frame}{Cache Performance and Locality}
\begin{columns}
\begin{column}{0.6\textwidth}
\concept{Locality Principles}:
\begin{itemize}
\item \textbf{Temporal Locality}: Recently accessed data likely to be accessed again
\item \textbf{Spatial Locality}: Nearby data likely to be accessed soon
\item \textbf{Sequential Locality}: Instructions executed in sequence
\item \textbf{Loop Locality}: Repeated access to same code/data
\end{itemize}

\vspace{0.3cm}
\concept{Cache Performance Metrics}:
\begin{itemize}
\item \textbf{Hit Rate}: Percentage of accesses found in cache
\item \textbf{Miss Rate}: Percentage of accesses not in cache
\item \textbf{Hit Time}: Time to access cache on hit
\item \textbf{Miss Penalty}: Time to handle cache miss
\end{itemize}

\vspace{0.3cm}
\concept{Harvard Cache Advantages}:
\begin{itemize}
\item No instruction-data conflicts in L1
\item Independent optimization of I-cache and D-cache
\item Better hit rates due to specialized caches
\item Simultaneous instruction fetch and data access
\end{itemize}

\vspace{0.3cm}
\concept{Performance Impact}:
\begin{itemize}
\item L1 miss resolved by L2 (much faster than main memory)
\item Multi-level hierarchy reduces average access time
\item Critical for modern high-performance processors
\end{itemize}
\end{column}
\begin{column}{0.4\textwidth}
\begin{center}
\textbf{Cache Performance}
\vspace{0.3cm}

\tikz{
% Performance graph
\draw[->] (0,0) -- (3,0);
\draw[->] (0,0) -- (0,3);
\node at (3.2,0) {Cache Level};
\node at (0,3.2) {Speed};

% Performance curve
\draw[thick, blue] (0.3,2.8) -- (1,2.2) -- (2,1.5) -- (2.8,0.5);

% Cache levels
\node at (0.3,-0.3) {L1};
\node at (1,-0.3) {L2};
\node at (2,-0.3) {L3};
\node at (2.8,-0.3) {Main};

% Access times
\node at (0.3,2.5) {1 cycle};
\node at (1,1.9) {10 cycles};
\node at (2,1.2) {30 cycles};
\node at (2.8,0.2) {300 cycles};
}

\vspace{0.5cm}

% Locality illustration
\tikz{
\node at (1.5,1.5) {\textbf{Locality Types}};

% Temporal
\draw[thick] (0,1) rectangle (1,1.2);
\node at (0.5,1.1) {Data};
\draw[->, red] (0.5,0.8) arc (180:0:0.3);
\node at (0.5,0.5) {Temporal};

% Spatial
\draw[thick] (1.5,1) rectangle (2.5,1.2);
\node at (2,1.1) {Array};
\draw[->, blue] (1.7,0.8) -- (2.3,0.8);
\node at (2,0.5) {Spatial};
}
\end{center}
\end{column}
\end{columns}
\end{frame}

% subsection 6: Implementation Considerations
\subsection{Implementation Considerations}

\begin{frame}{Modern Harvard Implementations}
\begin{columns}
\begin{column}{0.6\textwidth}
\concept{Contemporary Applications}:
\begin{itemize}
\item \textbf{Microcontrollers}: Embedded systems, IoT devices
\item \textbf{DSP Processors}: Digital signal processing applications
\item \textbf{GPU Architectures}: Parallel processing units
\item \textbf{RISC Processors}: High-performance computing systems
\end{itemize}

\vspace{0.3cm}
\concept{Implementation Challenges}:
\begin{itemize}
\item \textbf{Memory Bandwidth}: Dual memory systems require more bandwidth
\item \textbf{Synchronization}: Coordinating instruction and data access
\item \textbf{Cache Coherency}: Maintaining consistency across caches
\item \textbf{Power Consumption}: Multiple memory interfaces increase power
\end{itemize}

\vspace{0.3cm}
\concept{Design Optimizations}:
\begin{itemize}
\item \textbf{Memory Controllers}: Specialized controllers for each memory
\item \textbf{Prefetching}: Predictive loading of instructions and data
\item \textbf{Pipeline Design}: Optimized for simultaneous access patterns
\item \textbf{Power Management}: Dynamic frequency and voltage scaling
\end{itemize}

\vspace{0.3cm}
\concept{Future Trends}:
\begin{itemize}
\item Integration with AI/ML accelerators
\item Advanced cache hierarchies
\item Non-volatile memory technologies
\item Quantum computing interfaces
\end{itemize}
\end{column}
\begin{column}{0.4\textwidth}
\begin{center}
\textbf{Modern Implementation}
\vspace{0.3cm}

\tikz{
% Modern processor
\draw[thick] (0.5,3) rectangle (2.5,3.5);
\node at (1.5,3.25) {Modern CPU};

% Multiple cores
\foreach \x in {0.7, 1.1, 1.5, 1.9, 2.3} {
  \draw[thick] (\x,2.5) rectangle (\x+0.2,2.8);
  \node at (\x+0.1,2.65) {C};
}

% Cache hierarchy
\draw[thick, fill=green!20] (0.5,2.1) rectangle (2.5,2.4);
\node at (1.5,2.25) {L1 I/D Caches};

\draw[thick, fill=blue!20] (0.3,1.7) rectangle (2.7,2);
\node at (1.5,1.85) {L2 Cache};

\draw[thick, fill=yellow!20] (0.1,1.3) rectangle (2.9,1.6);
\node at (1.5,1.45) {L3 Cache};

% Memory controllers
\draw[thick] (0,0.8) rectangle (1.4,1.1);
\node at (0.7,0.95) {I-Mem Ctrl};

\draw[thick] (1.6,0.8) rectangle (3,1.1);
\node at (2.3,0.95) {D-Mem Ctrl};

% Connections
\draw[<->] (1.2,1.3) -- (0.7,1.1);
\draw[<->] (1.8,1.3) -- (2.3,1.1);

% Labels
\node at (1.5,0.5) {\textbf{Harvard Architecture}};
\node at (1.5,0.2) {Separate I/D Paths};
}
\end{center}
\end{column}
\end{columns}
\end{frame}

% Summary
\subsection{Summary}

\begin{frame}{Chapter Summary}
\begin{columns}
\begin{column}{0.6\textwidth}
\concept{Key Concepts Covered}:
\begin{itemize}
\item Fifth-Order Systems (5-OS) and Harvard abstract model
\item Historical development from Turing machines to modern computers
\item Comparison between Harvard and Von Neumann models
\item RISC principles and hardware-software segregation
\item Cache memory hierarchy and performance implications
\item Modern implementation considerations and challenges
\end{itemize}

\vspace{0.3cm}
\concept{Architectural Innovations}:
\begin{itemize}
\item Separate instruction and data memories
\item Simultaneous access capabilities
\item Cache hierarchy optimization
\item RISC instruction set design
\item Performance-oriented architecture
\end{itemize}

\vspace{0.3cm}
\concept{Modern Relevance}:
\begin{itemize}
\item Foundation for high-performance processors
\item Basis for embedded system design
\item Influence on GPU architectures
\item Principles in AI/ML accelerators
\item Continued evolution in computer architecture
\end{itemize}
\end{column}
\begin{column}{0.4\textwidth}
\begin{center}
\textbf{Harvard Model Evolution}
\vspace{0.3cm}

\tikz{
\node[draw, circle, minimum width=1.5cm] (center) at (0,0) {Harvard};
\node[draw, circle, minimum width=1.5cm] (center2) at (0,0) {5-OS};
\node[draw, rectangle] (hist) at (0,1.5) {Historical};
\node[draw, rectangle] (risc) at (-1.5,0.5) {RISC};
\node[draw, rectangle] (cache) at (1.5,0.5) {Cache};
\node[draw, rectangle] (modern) at (0,-1.5) {Modern};

\draw[->] (hist) -- (center);
\draw[->] (center) -- (risc);
\draw[->] (center) -- (cache);
\draw[->] (center) -- (modern);
}
\end{center}
\end{column}
\end{columns}
\end{frame}

\begin{frame}{Discussion Questions}
\begin{enumerate}
\item How does the Harvard model address the von Neumann bottleneck?
\item What are the trade-offs between Harvard and von Neumann architectures?
\item How do RISC principles complement the Harvard architecture?
\item What role does cache hierarchy play in modern Harvard implementations?
\item How has the Harvard model influenced modern processor design?
\item What are the implications of separate instruction and data memories?
\item How do modern applications benefit from Harvard architecture principles?
\item What future developments might further evolve the Harvard model?
\end{enumerate}
\end{frame}

